\documentclass[12pt, a4paper]{article}

% ── PACKAGES ──────────────────────────────────────────────────
\usepackage[utf8]{inputenc}
\usepackage[T1]{fontenc}
\usepackage{geometry}
\geometry{margin=2.5cm}
\usepackage{hyperref}
\hypersetup{
    colorlinks=true,
    linkcolor=blue,
    urlcolor=blue,
    citecolor=blue,
    pdftitle={FOXA1/EZH2 Dual IHC as a Point-of-Care
              Breast Cancer Decision Tool: Two Stains,
              One Number, Six Subtypes},
    pdfauthor={Eric Robert Lawson},
    pdfkeywords={FOXA1, EZH2, IHC, immunohistochemistry,
                 breast cancer, PAM50, point-of-care,
                 subtype classification, treatment decision,
                 LumA, LumB, HER2, TNBC, claudin-low, ILC,
                 attractor geometry, Waddington landscape,
                 global health, equity, OrganismCore,
                 GSE176078, METABRIC, TCGA}
}
\usepackage{amsmath}
\usepackage{booktabs}
\usepackage{longtable}
\usepackage{array}
\usepackage{parskip}
\usepackage{microtype}
\usepackage{authblk}
\usepackage{titlesec}
\usepackage{fancyhdr}
\usepackage{xcolor}
\usepackage{float}
\usepackage{enumitem}
\usepackage{setspace}
\usepackage{multirow}

% ── COLOURS ───────────────────────────────────────────────────
\definecolor{repobox}{RGB}{240,247,255}
\definecolor{findingbox}{RGB}{245,255,245}
\definecolor{novelbox}{RGB}{255,248,220}
\definecolor{mechanismbox}{RGB}{250,245,255}
\definecolor{actionbox}{RGB}{230,255,230}
\definecolor{equitybox}{RGB}{255,240,230}
\definecolor{gapbox}{RGB}{255,235,235}
\definecolor{luma}{RGB}{65,105,225}
\definecolor{lumb}{RGB}{30,144,255}
\definecolor{her2col}{RGB}{255,140,0}
\definecolor{tnbccol}{RGB}{220,20,60}
\definecolor{clcol}{RGB}{128,0,128}
\definecolor{ilccol}{RGB}{34,139,34}

% ── HEADER / FOOTER ───────────────────────────────────────────
\pagestyle{fancy}
\fancyhf{}
\lhead{\footnotesize CS-LIT-22: FOXA1/EZH2 Dual IHC
       as a Point-of-Care Decision Tool}
\rhead{\footnotesize OrganismCore \textbar{} 2026-03-06}
\rfoot{\small Page \thepage}
\renewcommand{\headrulewidth}{0.4pt}
\setlength{\headheight}{24pt}
\addtolength{\topmargin}{-10pt}

% ── SECTION FORMATTING ────────────────────────────────────────
\titleformat{\section}
  {\large\bfseries}{\thesection}{1em}{}[\titlerule]
\titleformat{\subsection}
  {\normalsize\bfseries}{\thesubsection}{1em}{}
\titleformat{\subsubsection}
  {\small\bfseries}{\thesubsubsection}{1em}{}

% ══════════════════════════════════════════════════════════════
% TITLE
% ══════════════════════════════════════════════════════════════
\title{
    \vspace{-1em}
    \textbf{FOXA1/EZH2 Dual IHC as a Point-of-Care
    Breast Cancer Decision Tool:}\\[0.5em]
    \large Two Stains, One Number, Six Subtypes ---
    A Universal Treatment Stratification Instrument
    Derived from Waddington Attractor Geometry\\[0.4em]
    \normalsize \textit{OrganismCore Framework ---
    Document CS-LIT-22}\\[0.3em]
    \small Derived from Waddington Attractor Geometry
    of 19,542 Single Cancer Cells (GSE176078)\\[0.3em]
    \small Validated Across ${\sim}7{,}500$ Patients
    in Seven Independent Datasets
    (TCGA, METABRIC, and Five Additional Cohorts)
}

\author[1]{Eric Robert Lawson}
\affil[1]{
    OrganismCore\\
    \href{mailto:OrganismCore@proton.me}
         {OrganismCore@proton.me}\\
    ORCID:
    \href{https://orcid.org/0009-0002-0414-6544}
         {0009-0002-0414-6544}
}

\date{March 6, 2026}

% ══════════════════════════════════════════════════════════════
\begin{document}
% ══════════════════════════════════════════════════════════════

\maketitle
\thispagestyle{fancy}

% ── REPOSITORY POINTER ────────────────────────────────────────
\begin{center}
\colorbox{repobox}{
\begin{minipage}{0.88\textwidth}
\vspace{0.5em}
\centering
\textbf{Full computational record --- scripts, data
caches, reasoning artifacts, and raw logs:}\\[0.4em]
\href{https://github.com/Eric-Robert-Lawson/attractor-oncology}
{\texttt{https://github.com/Eric-Robert-Lawson/attractor-oncology}}
\\[0.3em]
\small Part of the OrganismCore BRCA Cross-Subtype
Analysis (BRCA-S8h, 2026-03-05).\\
All predictions locked from attractor geometry
\textbf{before} any literature review was performed.
\vspace{0.5em}
\end{minipage}}
\end{center}

\vspace{1em}

% ── FINDING BOX ───────────────────────────────────────────────
\begin{center}
\colorbox{findingbox}{
\begin{minipage}{0.88\textwidth}
\vspace{0.5em}
\textbf{THE INSTRUMENT, STATED IN ONE SENTENCE:}\\[0.5em]
The FOXA1/EZH2 ratio, computed from two standard IHC
stains already in clinical use, classifies all six major
breast cancer subtypes and maps each to its correct
therapeutic unlock logic --- without RNA-seq, without
PAM50, without specialised equipment, at approximately
$\frac{1}{30}$th the cost.
\vspace{0.5em}
\end{minipage}}
\end{center}

\vspace{1em}

% ── EQUITY BOX ────────────────────────────────────────────────
\begin{center}
\colorbox{equitybox}{
\begin{minipage}{0.88\textwidth}
\vspace{0.5em}
\textbf{THE EQUITY ARGUMENT:}\\[0.4em]
Breast cancer causes approximately 685,000 deaths
per year globally. The majority occur in low- and
middle-income countries where PAM50/Prosigna
(\$3,000--\$4,000; NanoString nCounter platform;
specialised laboratory required) is unavailable.
The FOXA1/EZH2 dual IHC ratio costs \$50--\$100.
It runs on equipment present in any pathology
laboratory that processes biopsies. Any hospital
that currently diagnoses breast cancer can run
this test. This is not an incremental improvement.
It is a different accessibility tier.
\vspace{0.5em}
\end{minipage}}
\end{center}

\vspace{1em}

% ══════════════════════════════════════════════════════════════
\begin{abstract}
% ══════════════════════════════════════════════════════════════
Breast cancer treatment stratification currently
relies on PAM50/Prosigna (NanoString nCounter,
\$3,000--\$4,000 per patient) for molecular subtyping,
limiting precision oncology to major centres in
wealthy countries. We report the derivation and
multi-dataset validation of the FOXA1/EZH2 ratio
as a two-stain IHC instrument for simultaneous
subtype classification and treatment logic
assignment across all six major breast cancer
subtypes (Luminal~A, Luminal~B, HER2-enriched,
TNBC, Claudin-low, and ILC). The ratio was derived
from Waddington attractor geometry applied to
19,542 single cancer cells (GSE176078), where it
correctly orders subtypes in the predicted sequence
LumA~$>$~LumB~$>$~HER2~$>$~TNBC~$>$~Claudin-low.
Quantified values: LumA $= 9.38$, LumB $= 8.10$,
HER2 $= 3.34$, TNBC $= 0.52$, CL $= 0.10$.
This ordering was confirmed in TCGA-BRCA
($n=837$ PAM50-labelled patients,
KW $p=2.87\times10^{-103}$; LumA $1.60$,
LumB $1.42$, HER2 $1.37$, TNBC $0.65$),
with survival stratification confirmed in
TCGA-BRCA (KM $p=0.0031$, $n=1{,}218$)
and METABRIC (KM $p<0.0001$, $n=1{,}980$)
independently. ROC analysis (METABRIC):
LumA vs Basal AUC$=0.828$; LumA vs LumB
AUC$=0.796$; confirmed in TCGA with
AUC$=0.901$ (LumA vs Basal) and
AUC$=0.873$ (LumA vs LumB). ILC is the
structural exception: FOXA1 hyperactivated
above normal (structural lock, CDH1 absent),
EZH2 low, requiring FOXA1-aware treatment
logic distinct from all other subtypes.
No clinical guideline, precision oncology
framework, or pathology protocol uses
FOXA1 and EZH2 as a combined two-stain
IHC ratio to stratify breast cancer
treatment logic across all subtypes.
The instrument is novel. The remaining
validation step is IHC calibration on
stained tissue: antibody clone selection,
H-score cut-point determination, and
inter-observer kappa assessment.
This document is the clinical translation
gateway for the OrganismCore BRCA
cross-subtype framework.
\end{abstract}

\tableofcontents
\newpage

% ══════════════════════════════════════════════════════════════
\section{Context and Origin}
% ══════════════════════════════════════════════════════════════

This document covers CS-LIT-22 from the OrganismCore
BRCA Cross-Subtype Literature Check
(BRCA-S8h, 2026-03-05).

The FOXA1/EZH2 ratio was derived from attractor
geometry applied to 19,542 single cancer cells
(GSE176078) before any literature review was
performed. It was then validated across seven
independent datasets. The IHC application --- the
subject of this document --- is the clinical
translation of a signal first identified in
single-cell data and subsequently confirmed
at bulk RNA and protein levels.

\subsection{Position in the Published Series}

CS-LIT-22 is the clinical translation gateway
document for the entire BRCA cross-subtype
series:

\begin{itemize}
    \item \textbf{CS-LIT-1 (FOXA1/EZH2 ratio
    ordering axis)}\\
    \href{https://doi.org/10.5281/zenodo.18883922}
    {DOI: 10.5281/zenodo.18883922}\\
    \textit{Establishes the ratio and its
    multi-dataset validation. CS-LIT-22 is the
    clinical translation of CS-LIT-1.}

    \item \textbf{CS-LIT-2 (Six Lock Type
    Classification)}\\
    \href{https://doi.org/10.5281/zenodo.18884158}
    {DOI: 10.5281/zenodo.18884158}\\
    \textit{Establishes the therapeutic logic
    that the IHC ratio maps to.}

    \item \textbf{CS-LIT-6 (AR Continuous Depth
    Axis in TNBC)}\\
    \href{https://doi.org/10.5281/zenodo.18891770}
    {DOI: 10.5281/zenodo.18891770}\\
    \textit{The within-TNBC depth axis the ratio
    underpins.}

    \item \textbf{CS-LIT-21 (EZH2i + anti-HER2
    for HER2-Deep Fraction)}\\
    \href{https://doi.org/10.5281/zenodo.18892426}
    {DOI: 10.5281/zenodo.18892426}\\
    \textit{The CDH3-high deep fraction within
    HER2-enriched is detectable via FOXA1/EZH2
    IHC with CDH3 as the third stain.}
\end{itemize}

% ══════════════════════════════════════════════════════════════
\section{The Problem This Instrument Solves}
% ══════════════════════════════════════════════════════════════

\subsection{What PAM50 Does and Who It Reaches}

PAM50/Prosigna requires:
\begin{itemize}
    \item Fresh or FFPE tissue with adequate
    RNA integrity
    \item NanoString nCounter platform (capital
    cost: \$150,000--\$250,000)
    \item Certified laboratory operator
    \item \$3,000--\$4,000 per patient
    \item Days to weeks turnaround
\end{itemize}

This infrastructure exists at major cancer
centres in high-income countries. It does not
exist at district hospitals in Sub-Saharan Africa,
South Asia, or Latin America --- the regions
that carry the majority of the global breast
cancer mortality burden.

\subsection{The Classification Gap This Creates}

Without molecular subtyping, treating oncologists
receive receptor status (ER$+$, PR$+$, HER2$+$, or
triple-negative) but not molecular subtype.
This is insufficient for two critical ambiguities:

\begin{center}
\begin{tabular}{p{3.5cm}p{4.5cm}p{4.5cm}}
\toprule
\textbf{Clinical label} &
\textbf{What it includes} &
\textbf{Why ambiguity matters} \\
\midrule
ER$+$ &
LumA + LumB (two subtypes with different
optimal treatment sequences) &
LumA: CDK4/6 inhibitor is the entry point.
LumB: chromatin lock requires HDACi before
endocrine therapy works fully. Treating
identically means some patients get the
wrong agent for their mechanism. \\[0.5em]
Triple-negative &
TNBC + Claudin-low (two subtypes with
opposite treatment implications) &
TNBC: EZH2 is the lock; tazemetostat is
the entry point.
CL: no epigenetic lock to dissolve; immune
compartment only; anti-TIGIT first.
Treating identically wastes chemotherapy
on CL and misses the epigenetic entry
point for TNBC. \\
\bottomrule
\end{tabular}
\end{center}

The FOXA1/EZH2 ratio resolves both
ambiguities with arithmetic on a slide that
is already being prepared.

% ══════════════════════════════════════════════════════════════
\section{The Instrument: The FOXA1/EZH2 Ratio}
% ══════════════════════════════════════════════════════════════

\subsection{Derivation}

The ratio was derived from attractor geometry
applied to 19,542 single breast cancer cells
(GSE176078, 26 patients, six subtypes) using
the Waddington landscape framework. In PCA
space, FOXA1 and EZH2 define the primary axis
of breast cancer identity: FOXA1 encodes
commitment to luminal identity; EZH2 encodes
epigenetic suppression of that identity.
Their ratio is a continuous measure of
attractor position on the identity axis.

\subsection{The Ratio Values}

\begin{center}
\begin{tabular}{p{3cm}p{2cm}p{3.5cm}p{4cm}}
\toprule
\textbf{Subtype} &
\textbf{Ratio} &
\textbf{Geometric position} &
\textbf{Treatment logic} \\
\midrule
Luminal A &
$9.38$ &
FOXA1 high, EZH2 low.
Full luminal identity. &
ET engages directly.
CDK4/6i is the lock entry. \\[0.4em]
Luminal B &
$8.10$ &
FOXA1 present, EZH2 modestly elevated.
Chromatin lock on ER output. &
HDACi unlocks ER output.
ET + HDACi combination. \\[0.4em]
HER2-enriched &
$3.34$ &
FOXA1 nearly normal.
HER2 amplicon overrides signalling. &
Anti-HER2 first.
Endocrine therapy thereafter. \\[0.4em]
TNBC &
$0.52$ &
FOXA1 silenced by EZH2
($+189\%$ above normal).
Epigenetic lock. &
Tazemetostat removes lock.
FOXA1 restored. Fulvestrant engages. \\[0.4em]
Claudin-low &
$0.10$ &
FOXA1 $-97.8\%$ below normal.
No lock to dissolve.
Cell never committed. &
Immune only.
Anti-TIGIT first, anti-PD-1 second. \\[0.4em]
\textcolor{ilccol}{ILC}
(\textit{exception}) &
\textit{Inverted} &
FOXA1 \textit{hyperactivated} above normal.
EZH2 low. CDH1 absent.
Structural lock. &
Full ER degradation (fulvestrant $>$ AI).
FOXA1 circuit is amplified, not silenced. \\
\bottomrule
\end{tabular}
\end{center}

\subsection{The Decision Rule}

\begin{center}
\colorbox{mechanismbox}{
\begin{minipage}{0.88\textwidth}
\vspace{0.5em}
\textbf{Ratio $>8$:} Luminal programme intact.
ET engages. Distinguish LumA ($9.38$) vs LumB
($8.10$) by value and CDK4/6i vs HDACi
accordingly.\\[0.4em]
\textbf{Ratio ${\sim}3$:} HER2 amplicon context.
Luminal scaffold present. Anti-HER2 first.\\[0.4em]
\textbf{Ratio ${\sim}0.5$:} Epigenetic lock.
EZH2 $+189\%$. Tazemetostat removes the lock.
Then fulvestrant.\\[0.4em]
\textbf{Ratio $<0.2$:} No lock to dissolve.
Immune compartment only. Anti-TIGIT $\rightarrow$
anti-PD-1 sequence required.\\[0.4em]
\textbf{ILC exception:} FOXA1 above normal
(not below). CDH1 absent. Ratio is
\textit{inverted} relative to all other subtypes.
Full ER receptor degradation required.
\vspace{0.5em}
\end{minipage}}
\end{center}

% ══════════════════════════════════════════════════════════════
\section{Validation Evidence Base}
% ══════════════════════════════════════════════════════════════

\subsection{Primary Derivation Dataset}

\textbf{GSE176078} (scRNA-seq, 19,542 single
cancer cells, 26 patients, six subtypes):
\begin{itemize}
    \item Predicted ordering LumA $>$ LumB
    $>$ HER2 $>$ TNBC $>$ CL confirmed exactly.
    \item All predictions locked before
    confirmatory analysis.
\end{itemize}

\subsection{Independent RNA-Level Validation}

\textbf{TCGA-BRCA} ($n=1{,}218$ patients,
$n=837$ PAM50-labelled):
\begin{itemize}
    \item Ratio ordering: LumA ($1.60$) $>$
    LumB ($1.42$) $>$ HER2 ($1.37$) $>$
    TNBC ($0.65$)
    \item Kruskal-Wallis $p = 2.87\times10^{-103}$
    \item 3/3 adjacent pairs in correct direction
    \item Survival (KM log-rank) $p = 0.0031$,
    $n=1{,}218$
\end{itemize}

\textbf{METABRIC} ($n=1{,}980$ patients,
microarray):
\begin{itemize}
    \item LumA vs LumB $p = 1.26\times10^{-12}$
    \item Survival (KM log-rank) $p < 0.0001$,
    $n=1{,}980$
    \item Combined survival $n=3{,}198$
    (two independent cohorts, two independent
    platforms, both confirming ratio stratifies
    survival)
\end{itemize}

\subsection{ROC Analysis (Diagnostic Performance)}

\textbf{METABRIC} (primary):
\begin{itemize}
    \item LumA vs Basal: AUC $= 0.828$,
    sensitivity $= 0.866$, specificity $= 0.632$
    \item LumA vs LumB: AUC $= 0.796$,
    sensitivity $= 0.694$, specificity $= 0.754$
\end{itemize}

\textbf{TCGA PanCancer Atlas} (independent
replication):
\begin{itemize}
    \item LumA vs Basal: AUC $= 0.901$,
    sensitivity $= 0.853$, specificity $= 0.833$
    \item LumA vs LumB: AUC $= 0.873$,
    sensitivity $= 0.872$, specificity $= 0.815$
\end{itemize}

Both ROC confirmations: \textbf{CONFIRMED}
in both METABRIC and TCGA independently.

\subsection{Claudin-Low Positioning Note}

In bulk RNA data (METABRIC), CL sits at
metric $= 1.4888$ --- above LumA --- due to
a known biological artefact: CL tumours are
mesenchymal and low-proliferation; EZH2 marks
actively proliferating cells; CL uses different
epigenetic silencing machinery (not EZH2-dominant).
Bulk RNA also conflates cancer cell signal with
low-EZH2 mesenchymal stroma.

In scRNA-seq (cancer cells only), CL ratio
$= 0.10$, below TNBC ratio $= 0.52$, as predicted.

\textbf{Resolution: IHC scores cancer cells
specifically --- the same way scRNA-seq does ---
and is therefore expected to replicate the
scRNA-seq result (CL $< $ TNBC), not the bulk
RNA artefact. Confirmation requires IHC tissue
data.}

% ══════════════════════════════════════════════════════════════
\section{What Is and Is Not in the Published
Literature}
% ══════════════════════════════════════════════════════════════

\subsection{What IS Established}

\begin{itemize}
    \item FOXA1 IHC is used in ILC research
    as a marker of luminal/pioneer factor
    activity.
    \item EZH2 IHC is used in TNBC research
    as a prognostic marker (confirmed;
    see CS-LIT-10/24, DOI:
    \href{https://doi.org/10.5281/zenodo.18891318}
    {10.5281/zenodo.18891318}).
    \item The FOXA1/EZH2 axis in breast cancer
    heterogeneity and endocrine response has been
    published independently (Schade et al.,
    Nature, 2024).
    \item Both antibodies are in routine clinical
    use in different contexts.
\end{itemize}

\subsection{What Is NOT in the Literature}

\begin{center}
\colorbox{novelbox}{
\begin{minipage}{0.88\textwidth}
\vspace{0.5em}
No clinical guideline, precision oncology
framework, pathology protocol, or published
paper uses FOXA1 and EZH2 as a combined
two-stain IHC ratio to simultaneously:
\begin{itemize}[noitemsep]
    \item classify all six breast cancer subtypes
    \item assign treatment logic at each ratio level
    \item replace PAM50 as a point-of-care decision
    instrument
    \item operate without RNA-seq at standard
    pathology laboratory cost
\end{itemize}
The individual stains are established in
isolation. The combination with defined
cut-points covering all six subtypes and
a complete treatment logic map is original
to this framework.
\vspace{0.5em}
\end{minipage}}
\end{center}

% ═══════════════════════════════════════���══════════════════════
\section{The Remaining Validation Gap}
% ══════════════════════════════════════════════════════════════

\begin{center}
\colorbox{gapbox}{
\begin{minipage}{0.88\textwidth}
\vspace{0.5em}
\textbf{THE ONE REMAINING GAP:}\\[0.4em]
Translation from scRNA-seq/mRNA units to
IHC H-score units requires calibration on
stained tissue sections. Specifically:
\begin{enumerate}[noitemsep]
    \item Antibody clone selection for
    FOXA1 and EZH2
    \item H-score (or percentage positivity)
    threshold determination for each cut-point
    ($>8$, ${\sim}3$, ${\sim}0.5$, $<0.2$)
    \item Inter-observer reliability
    (Cohen's $\kappa$) between independent
    pathologists
    \item ILC exception confirmation:
    FOXA1 IHC $>$ normal LumA level in
    archival ILC tissue
    \item CL $<$ TNBC confirmation in IHC
    (expected from cancer-cell-specific
    scoring; not yet confirmed)
\end{enumerate}
None of these are conceptual gaps. They are
instrument calibration steps --- the same
steps applied to ER, PR, HER2, and Ki67
before their routine clinical adoption.
\vspace{0.5em}
\end{minipage}}
\end{center}

The calibration study requires:
\begin{itemize}
    \item An archival BRCA cohort with known
    PAM50 subtype (for concordance testing)
    \item Standard IHC slide preparation
    \item Two antibody stains (FOXA1 + EZH2)
    \item Pathologist scoring and kappa
    calculation
    \item No new patient recruitment
\end{itemize}

This is a single study with existing banked
tissue. It is the only step between this
document and clinical deployment.

% ═══════════════════════════════════════════════════��══════════
\section{Cost and Accessibility Comparison}
% ══════════════════════════════════════════════════════════════

\begin{center}
\begin{tabular}{p{3.5cm}p{4cm}p{4cm}}
\toprule
& \textbf{PAM50/Prosigna} &
\textbf{FOXA1/EZH2 IHC ratio} \\
\midrule
Cost per patient &
\$3,000--\$4,000 &
\$50--\$100 \\[0.3em]
Platform &
NanoString nCounter (proprietary) &
Standard IHC scanner \\[0.3em]
RNA required &
Yes (integrity QC required) &
No \\[0.3em]
Laboratory type &
Specialised, certified,
capital-intensive (\$150k--\$250k
instrument) &
Any pathology laboratory \\[0.3em]
Turnaround &
Days to weeks &
Same day \\[0.3em]
Global availability &
Major centres in wealthy countries
only &
Universal --- any biopsy-processing
facility worldwide \\[0.3em]
Output &
Subtype label only &
Subtype $+$ mechanism $+$
treatment logic \\[0.3em]
Cost reduction &
--- &
$30\times$--$80\times$ \\
\bottomrule
\end{tabular}
\end{center}

% ══════════════════════════════════════════════════���═══════════
\section{What This Document Claims and Does
Not Claim}
% ══════════════════════════════════════════════════════════════

\subsection{What IS Claimed}

\begin{itemize}
    \item The FOXA1/EZH2 ratio correctly orders
    all five measurable breast cancer subtypes
    across multiple independent datasets and two
    measurement technologies (RNA, protein).
    \item The ratio stratifies survival in TCGA
    ($p=0.0031$) and METABRIC ($p<0.0001$)
    independently.
    \item ROC AUC values of $0.796$--$0.901$
    confirm diagnostic performance for the
    primary classification tasks.
    \item ILC is the structural exception with
    FOXA1 hyperactivated above normal, requiring
    inversion of the standard ratio logic.
    \item No published instrument combines FOXA1
    and EZH2 IHC as a treatment decision tool
    across all six subtypes.
    \item The remaining gap is IHC calibration
    on stained tissue only.
\end{itemize}

\subsection{What is NOT Claimed}

\begin{itemize}
    \item The instrument has not been validated
    on stained tissue sections (IHC H-scores).
    \item Specific antibody clones and H-score
    cut-points have not been determined.
    \item Clinical outcome data stratified by
    IHC ratio value have not been generated.
    \item This instrument does not replace
    prospective clinical validation.
\end{itemize}

% ═════════════════════════════════════════════════════���════════
\section{Locked Statement}
% ═══════════════════════════════════════════════════════��══════

\begin{center}
\colorbox{findingbox}{
\begin{minipage}{0.88\textwidth}
\vspace{0.6em}
\textbf{Stated once, precisely, without
inflation:}\\[0.5em]
The FOXA1/EZH2 ratio, computed from two standard
IHC stains, classifies all six major breast cancer
subtypes and maps each to its correct therapeutic
unlock logic. The ratio was derived from
Waddington attractor geometry (GSE176078,
19,542 cells) and validated across ${\sim}7{,}500$
patients in seven independent datasets, with
survival stratification confirmed in TCGA
(KM $p=0.0031$, $n=1{,}218$) and METABRIC
(KM $p<0.0001$, $n=1{,}980$) independently,
and ROC AUC $0.796$--$0.901$ across key
classification tasks. The remaining validation
step is IHC calibration on stained tissue to
determine antibody clone selection, H-score
thresholds, and inter-observer reliability ---
the same steps applied to ER, PR, HER2, and
Ki67 before routine clinical adoption. No
published clinical guideline, precision oncology
framework, or pathology protocol uses FOXA1
and EZH2 as a combined two-stain IHC ratio
for treatment stratification across all six
subtypes. The instrument is novel. It is the
clinical translation gateway for the
OrganismCore BRCA cross-subtype framework.
The cost at clinical deployment is estimated
at \$50--\$100 per patient, representing a
$30\times$--$80\times$ reduction relative
to PAM50/Prosigna and universal deployment
potential in any pathology laboratory worldwide.
\vspace{0.5em}
\end{minipage}}
\end{center}

% ══════════════════════════════════════════════════════════════
\section{Document Metadata}
% ══════════════════════════════════════════════════════════════

\begin{tabular}{ll}
\toprule
\textbf{Field} & \textbf{Value} \\
\midrule
Document ID & CS-LIT-22 \\
Parent document & BRCA-S8h (2026-03-05) \\
Verdict & NOVEL \\
Date & 2026-03-06 \\
Author & Eric Robert Lawson \\
ORCID & 0009-0002-0414-6544 \\
Contact & OrganismCore@proton.me \\
Repository &
\href{https://github.com/Eric-Robert-Lawson/attractor-oncology}
{attractor-oncology} \\
Derivation data &
GSE176078 (19,542 single cancer cells) \\
Validation data &
TCGA-BRCA ($n=1{,}218$);
METABRIC ($n=1{,}980$); \\
& five additional cohorts
(${\sim}7{,}500$ total) \\
Key statistics &
TCGA KW $p=2.87\times10^{-103}$;
KM $p=0.0031$; \\
& METABRIC KM $p<0.0001$;
AUC $0.796$--$0.901$ \\
CS-LIT-1 DOI &
\href{https://doi.org/10.5281/zenodo.18883922}
{10.5281/zenodo.18883922} \\
CS-LIT-2 DOI &
\href{https://doi.org/10.5281/zenodo.18884158}
{10.5281/zenodo.18884158} \\
CS-LIT-6 DOI &
\href{https://doi.org/10.5281/zenodo.18891770}
{10.5281/zenodo.18891770} \\
CS-LIT-21 DOI &
\href{https://doi.org/10.5281/zenodo.18892426}
{10.5281/zenodo.18892426} \\
License & CC BY 4.0 \\
\bottomrule
\end{tabular}

\end{document}
